\chapter{Analityczne podstawy teorii Seiberga--Wittena: przestrzenie moduli i komory}
\label{ch:tg-sw-analiza}

Poprzedni rozdział zdefiniował równania Seiberga--Wittena, wskazał ich
wymiar oczekiwany i zastosowania do wymiernej kontrakcji. Aby zrozumieć,
dlaczego liczenie rozwiązań daje niezmiennik różniczkowy, potrzebujemy
jeszcze aparatu analitycznego. Omówimy kolejno dopełnienia Sobolewa,
iloraz przez grupę cechowania, kompleks deformacyjny, zwartość,
transwersalność i orientacje. Następnie wyjaśnimy, dlaczego przypadek
$b_2^+=1$ wymaga teorii komór. Podamy pełne argumenty algebraiczne
i zasadnicze oszacowania; dwa głębokie kroki analityczne, tj. globalną
teorię przekrojów Fredholma i lokalną teorię Kuranishiego w pobliżu
rozwiązań redukowalnych, sformułujemy jawnie jako przyjmowane twierdzenia.

\section{Dopełnienia Sobolewa i grupa cechowania}
\label{sec:tg-sw-sobolev}

Przez cały rozdział $X$ oznacza zamkniętą, spójną, zorientowaną,
gładką czterorozmaitość z ustaloną metryką riemannowską $g$,
a $\mathfrak s$ --- strukturę $\mathrm{Spin}^c$ z wiązkami $S^\pm$
i wiązką wyznacznikową $L$. Oznaczamy przez
$\Omega^j(X;i\mathbb R)$ przestrzeń gładkich $j$-form o wartościach
w $i\mathbb R$ oraz przez $\Gamma(S^\pm)$ przestrzenie spinorów.
Wybieramy raz na zawsze gładkie połączenie unitarne $A_0$ na $L$.

\begin{definicja}[Przestrzenie Sobolewa i konfiguracje]
Dla całkowitego $k\geq3$ niech $W^{k,2}(E)$ będzie dopełnieniem
przestrzeni gładkich przekrojów wiązki euklidesowej lub hermitowskiej
$E\to X$ w normie Sobolewa rzędu $k$, zdefiniowanej przez
wybraną koneksję i metrykę. \emph{Przestrzenią konfiguracji} nazywamy
przestrzeń afiniczną
\[
\mathcal C_k=(A_0+W^{k,2}(iT^*X))\times W^{k,2}(S^+).
\]
\emph{Grupą cechowania} jest $\mathcal G_{k+1}:=W^{k+1,2}(X,S^1)$,
z działaniem $u\cdot(A,\psi)=(A-2u^{-1}\mathrm du,u\psi)$.
Konfigurację nazywamy \emph{nieredukowalną}, jeżeli
$\psi\not\equiv0$; zbiór takich konfiguracji oznaczamy
$\mathcal C_k^*$.
\end{definicja}

W wymiarze cztery zachodzi $k>4/2$, dlatego
twierdzenie Sobolewa daje ciągłe zanurzenie
$W^{k,2}(X)\hookrightarrow C^0(X)$ i ciągłość
mnożenia $W^{k,2}\times W^{k,2}\to W^{k,2}$.
Analogiczne oszacowania dla przekrojów otrzymujemy
w lokalnych trywializacjach przy użyciu podziału jedności.
Grupa $\mathcal G_{k+1}$ jest grupą Liego
modelowaną na przestrzeni Hilberta:
lokalną mapę w otoczeniu $u_0$ daje
$f\mapsto u_0 e^{if}$ dla małych
$f\in W^{k+1,2}(X;\mathbb R)$.
Ciągłe zanurzenie w $C^0$ pozwala wybrać
jednoznaczną gałąź logarytmu w tym otoczeniu.
Gładkość mnożenia, odwracania i wykładnika wynika
z własności algebry Sobolewa.
Ponieważ $u^{-1}=\overline u$ oraz
$\mathrm du\in W^{k,2}$, działanie cechowania
jest gładkie na $\mathcal C_k$.
Do operatora równań używamy przestrzeni docelowej
$W^{k-1,2}(S^-\oplus i\Lambda^2_+)$:
operator Diraca i krzywizna zawierają pochodne
pierwszego rzędu, natomiast iloczyny połączenia
ze spinorem oraz $q(\psi)$ należą do $W^{k,2}$.
W szczególności operator równań jest gładki
między wskazanymi przestrzeniami Hilberta. W rachunkach
linearyzacyjnych należy odróżnić parametr rzeczywisty $f$,
odpowiadający jednoparametrowej grupie $u_t=e^{itf}$, od form
i-funkcji o wartościach w $i\mathbb R$.

\begin{stwierdzenie}[Stabilizator konfiguracji]
\label{prop:tg-sw-stabilizer}
Niech $(A,\psi)\in\mathcal C_k$. Jeżeli $u\in\mathcal G_{k+1}$
zachowuje tę konfigurację, to $u$ jest stałą funkcją o wartości
w $S^1$. Jeżeli ponadto $\psi\not\equiv0$, to $u=1$.
Wobec tego działanie $\mathcal G_{k+1}$ na $\mathcal C_k^*$
jest wolne, a stabilizatorem konfiguracji redukowalnej jest $S^1$.
\end{stwierdzenie}

\begin{proof}
Z równości $A-2u^{-1}\mathrm du=A$ wynika
$u^{-1}\mathrm du=0$ w sensie słabych pochodnych.
Mnożąc przez $u$ dostajemy $\mathrm du=0$.
Każda z rzeczywistych składowych $u$ jest
funkcją Sobolewa o zerowej słabej pochodnej,
zatem jest stała prawie wszędzie na każdym
spójnym otoczeniu współrzędnościowym.
Ciągłość reprezentanta Sobolewa daje równość
wszędzie. Funkcja $u$ jest więc lokalnie stała,
a ponieważ $X$ jest spójna, jest stała globalnie.
Równość $u\psi=\psi$ przy $\psi\not\equiv0$ wymusza
$u=1$ w punkcie, w którym $\psi$ jest niezerowy,
a więc na całym $X$. Przy $\psi=0$ każda stała faza
zachowuje konfigurację.
\end{proof}

\begin{uwaga}[Dlaczego potrzebny jest wybór przekroju?]
Iloraz $\mathcal C_k/\mathcal G_{k+1}$ nie jest na ogół gładką
rozmaitością w pobliżu konfiguracji redukowalnych, ponieważ
stabilizator nie jest tam trywialny. Na części nieredukowalnej
twierdzenie o przekroju lokalnym (ang. \emph{slice theorem})
daje lokalny model ilorazu. Warunek $\mathrm d^*a=0$
na różnicy połączeń $A-A_0=ia$ opisuje cechowanie Coulomba
samej części połączeniowej, ale nie usuwa resztkowego działania
stałych faz na spinor. Pełny liniowy warunek przekroju
w punkcie $(A,\psi)$ ma postać $\mathrm d_0^*(ia,\phi)=0$,
a zatem uwzględnia również kierunek spinorowy.
Nie należy utożsamiać takiego lokalnego przekroju
z globalnym wyborem cechowania na całej przestrzeni konfiguracji.
\end{uwaga}

\section{Przekrój Coulomba i kompleks deformacyjny}
\label{sec:tg-sw-deform}

Wybierzmy rozwiązanie $(A,\psi)$ równań z poprzedniego
rozdziału. Wygodnie zapisać ich lewą stronę jako
$\mathfrak F_\eta(A,\psi):=(D_A\psi,F_A^+-q(\psi)-i\eta)$.
Niech $\rho$ oznacza mnożenie Clifforda, a
$\mathrm dq_\psi$ różniczkę odwzorowania kwadratowego $q$
w punkcie $\psi$.

\begin{stwierdzenie}[Linearyzacja i różniczka działania]
\label{prop:tg-sw-linear}
W punkcie będącym rozwiązaniem równania różniczka działania
cechowania w kierunku $f\in\Omega^0(X;\mathbb R)$ oraz
różniczka operatora $\mathfrak F_\eta$ mają postać
\[
\begin{aligned}
\mathrm d_0f&=(-2i\,\mathrm df,\,if\psi),\\
\mathrm d_1(a,\phi)&=\bigl(D_A\phi+\tfrac12\rho(a)\psi,\,
\mathrm d^+a-(\mathrm dq)_\psi\phi\bigr),
\end{aligned}
\]
gdzie $a\in\Omega^1(X;i\mathbb R)$ i $\phi\in\Gamma(S^+)$.
Zachodzi $\mathrm d_1\circ\mathrm d_0=0$.
\end{stwierdzenie}

\begin{proof}
Dla $u_t=e^{itf}$ mamy
$u_t^{-1}\mathrm du_t=it\,\mathrm df$; różniczkując działanie
w $t=0$, otrzymujemy wzór na $\mathrm d_0$.
Krzywizna połączenia $A+ta$ jest równa
$F_A+t\,\mathrm da$, gdyż grupa strukturalna $U(1)$
jest przemienna. Stąd różniczka równania krzywizny
to $\mathrm d^+a-(\mathrm dq)_\psi\phi$.
Dla przyjętej konwencji połączenia na wiązce wyznacznikowej
różniczka $D_{A+ta}$ wynosi $\rho(a)/2$, co daje
pierwszą składową $\mathrm d_1$.
Bezpośrednio sprawdźmy także warunek kompleksu.
Reguła Leibniza operatora Diraca daje
$D_A(if\psi)=i\rho(\mathrm df)\psi+ifD_A\psi$.
Po podstawieniu $a=-2i\,\mathrm df$ pierwszy
składnik $\mathrm d_1\mathrm d_0f$ wynosi zatem
$ifD_A\psi=0$. Drugi znika, gdyż
$\mathrm d^+\mathrm df=0$, a
$(\mathrm dq)_\psi(if\psi)=0$ wynika
z różniczkowania punktowej równości
$q(e^{itf}\psi)=q(\psi)$.
Otrzymujemy $\mathrm d_1\mathrm d_0=0$.
Równoważnie można skorzystać z niezmienniczości
zbioru rozwiązań względem cechowania.
\end{proof}

\begin{lemat}[Warunek przekroju i stałe transformacje cechowania]
\label{lem:tg-sw-full-coulomb}
Niech $(A,\psi)$ będzie gładką konfiguracją nieredukowalną.
Na przestrzeni stycznej wyposażonej w rzeczywisty iloczyn
skalarny $L^2$ formalne sprzężenie operatora
$\mathrm d_0f=(-2i\,\mathrm df,if\psi)$ jest dane wzorem
\[
\mathrm d_0^*(i\alpha,\phi)
=-2\,\mathrm d^*\alpha
+\operatorname{Re}\langle i\psi,\phi\rangle,
\]
gdzie $\alpha\in\Omega^1(X;\mathbb R)$, a iloczyn
spinorów jest hermitowski. Ponadto
$\mathrm d_0^*\mathrm d_0=4\,\mathrm d^*\mathrm d+|\psi|^2$
na funkcjach rzeczywistych i operator ten ma zerowe jądro.
W konsekwencji w otoczeniu konfiguracji nieredukowalnej
równanie $\mathrm d_0^*(i\alpha,\phi)=0$ jest liniowym
warunkiem lokalnego przekroju dla działania cechowania.
\end{lemat}

\begin{proof}
Z definicji sprzężenia i dla dowolnej rzeczywistej funkcji
$f$ obliczamy
\[
\begin{aligned}
\langle \mathrm d_0f,(i\alpha,\phi)\rangle_{L^2}
&=-2\langle\mathrm df,\alpha\rangle_{L^2}
+\int_X f\,\operatorname{Re}\langle i\psi,\phi\rangle
\,\mathrm d\mu_g\\
&=\int_X f\bigl(-2\,\mathrm d^*\alpha
+\operatorname{Re}\langle i\psi,\phi\rangle\bigr)
\,\mathrm d\mu_g.
\end{aligned}
\]
Stąd otrzymujemy wzór na sprzężenie oraz
$\mathrm d_0^*\mathrm d_0f=4\,\mathrm d^*\mathrm df+|\psi|^2f$.
Jeśli ten operator zeruje $f$, to przez całkowanie
$4\|\mathrm df\|_{L^2}^2+\|f\psi\|_{L^2}^2=0$.
Zatem $f$ jest stała i $f\psi=0$. Ponieważ
$\psi\not\equiv0$, musi być $f=0$.
Eliptyczny, samosprzężony operator $\mathrm d_0^*\mathrm d_0$
ma więc trywialne jądro i, na odpowiednich dopełnieniach
Sobolewa, jest odwracalny. Istotnie, samosprzężoność i eliptyczność
na zamkniętej rozmaitości dają zerowy
indeks, a regularność eliptyczna zapewnia
izomorfizm $W^{k+1,2}\to W^{k-1,2}$.
Jeśli $\mathfrak c=(A,\psi)$, rozważamy
\[
(f,\xi)\longmapsto
\mathrm d_0^*(e^{if}\cdot(\mathfrak c+\xi)-\mathfrak c)
\]
na otoczeniu $(0,0)$ w iloczynie
$W^{k+1,2}(X;\mathbb R)$ i przestrzeni
stycznej klasy $W^{k,2}$.
Odwzorowanie jest gładkie, znika w $(0,0)$,
a jego różniczka względem $f$ jest
$\mathrm d_0^*\mathrm d_0$.
Twierdzenie o funkcji uwikłanej daje jedyną
małą transformację ustawiającą pobliską
konfigurację w przekroju. Szczegółowy opis
lokalnego modelu podajemy w twierdzeniu
\ref{thm:tg-sw12-slice}. Sama ta jednoznaczność
dotyczy otoczenia identyczności, a nie
wszystkich składowych grupy cechowania.
\end{proof}

\begin{definicja}[Kompleks deformacyjny]
\emph{Kompleksem deformacyjnym Seiberga--Wittena} w punkcie
$(A,\psi)$ będącym rozwiązaniem nazywamy kompleks rzeczywisty
\[
0\longrightarrow\Omega^0(X;\mathbb R)
\xrightarrow{\mathrm d_0}\Omega^1(X;i\mathbb R)\oplus\Gamma(S^+)
\xrightarrow{\mathrm d_1}\Gamma(S^-)\oplus\Omega^2_+(X;i\mathbb R)
\longrightarrow0.
\]
Pierwszy wyraz zapisujemy jako funkcje rzeczywiste,
ponieważ przyjęliśmy $u_t=e^{itf}$ oraz
$\mathrm d_0f=(-2i\,\mathrm df,if\psi)$.
Równoważna konwencja z $if\in\Omega^0(X;i\mathbb R)$
prowadzi do tej samej teorii po odpowiedniej zmianie wzoru
na $\mathrm d_0$. Jego kohomologie oznaczamy przez
$H^0_{A,\psi}$, $H^1_{A,\psi}$ i $H^2_{A,\psi}$.
\end{definicja}

\begin{uwaga}[Rola poszczególnych kohomologii]
$H^0_{A,\psi}$ jest algebrą Liego stabilizatora,
$H^1_{A,\psi}$ jest przestrzenią styczną do przestrzeni
moduli w przypadku regularnym, a $H^2_{A,\psi}$ jest
przestrzenią przeszkód w zastosowaniu twierdzenia o funkcji
uwikłanej. Dla rozwiązania nieredukowalnego
$H^0_{A,\psi}=0$. Gdy ponadto $H^2_{A,\psi}=0$, rozwiązanie
jest \emph{regularne} i przestrzeń moduli jest lokalnie gładka.
W przypadku nieregularnym $H^1_{A,\psi}$ opisuje
przestrzeń deformacji pierwszego rzędu, lecz nie musi być
przestrzenią styczną do gładkiej rozmaitości.
Są to stwierdzenia lokalne, po wykonaniu dopełnień Sobolewa
i użyciu twierdzenia o przekroju.
\end{uwaga}

\begin{twierdzenie}[Eliptyczność i indeks]
\label{thm:tg-sw-elliptic}
Kompleks deformacyjny jest eliptyczny, a dla $\mathfrak s$
o klasie $c=c_1(L)$ jego indeks rzeczywisty
$\dim H^1-\dim H^0-\dim H^2$ wynosi
\[
d(\mathfrak s)=\frac{\langle c^2,[X]\rangle-2e(X)-3\sigma(X)}4.
\]
\end{twierdzenie}

\begin{proof}[Wskazanie części elementarnej i twierdzenia indeksowego]
Symbole członów liniowych nie zależą od $\psi$:
człony z mnożeniem przez $\psi$ i różniczką $q$ mają rząd zero.
Główna część rozpada się więc na symbol kompleksu
$\Omega^0\xrightarrow{\mathrm d}\Omega^1
\xrightarrow{\mathrm d^+}\Omega^2_+$ oraz symbol operatora Diraca.
Dla niezerowego kowektora $\xi$ pomijamy
nieistotne odwracalne czynniki skalarne w symbolach.
Pierwsza mapa ma postać $f\mapsto f\xi$,
a druga $a\mapsto(\xi\wedge a)^+$.
Pierwsza jest injektywna. Rozkładamy
$a=t\xi+a_\perp$, gdzie $a_\perp\perp\xi$.
Forma $\beta=\xi\wedge a_\perp$ spełnia
$\beta\wedge\beta=0$, stąd
$\|\beta^+\|^2=\|\beta^-\|^2=\|\beta\|^2/2$.
Zatem $(\xi\wedge a)^+=0$ dokładnie wtedy,
gdy $a_\perp=0$. Otrzymujemy dokładność
w środkowym wyrazie. Ponieważ
$\dim_\mathbb R\xi^\perp=\dim_\mathbb R\Lambda^2_+=3$,
druga mapa jest surjektywna.
Symbol Diraca jest izomorfizmem, gdyż
$\rho(\xi)^*\rho(\xi)=|\xi|^2\operatorname{id}$.
To dowodzi dokładności całego ciągu symboli.
Na dopełnieniach klas $W^{k+1,2}$, $W^{k,2}$,
$W^{k-1,2}$ eliptyczny kompleks ma domknięte
obrazy i skończeniewymiarowe kohomologie.
Operator $\mathrm d_0^*\oplus\mathrm d_1$
jest Fredholmowski; jego jądro identyfikuje się
z $H^1$, a cokernel z $H^0\oplus H^2$,
na mocy rozkładu Hodge'a dla kompleksów eliptycznych.
Człony rzędu zero są zwarte między odpowiednimi
przestrzeniami Sobolewa, więc nie zmieniają
indeksu operatora.
Indeks części de Rhama wynosi $b_1-1-b_2^+$,
a indeks rzeczywisty Diraca to
$2\operatorname{ind}_{\mathbb C}D_A=(c^2-\sigma)/4$.
Po ich dodaniu i podstawieniu
$e=2-2b_1+b_2^++b_2^-$ oraz
$\sigma=b_2^+-b_2^-$ dostajemy wzór z tezy.
Wartość indeksu Diraca jest tutaj zastosowaniem
twierdzenia Atiyaha--Singera, którego nie dowodzimy.
\end{proof}

\section{Oszacowanie spinora i zwartość}
\label{sec:tg-sw-compact}

Kluczowa różnica między teorią Seiberga--Wittena a niektórymi
innymi teoriami cechowania polega na tym, że na zamkniętej
czterorozmaitości przestrzeń rozwiązań modulo cechowanie
jest zwarta bez dodawania punktów odpowiadających
koncentracji energii. Wyprowadzimy najważniejsze oszacowanie.

\begin{lemat}[Oszacowanie z formuły Weitzenböcka]
\label{lem:tg-sw-bound}
Przy ustalonej gładkiej metryce $g$, gładkiej perturbacji
$\eta\in\Omega^2_+(X)$ i normalizacji odwzorowania $q$
jak w rozdziale poprzednim istnieją stałe $c_0,C_0>0$,
zależne tylko od przyjętych konwencji algebraicznych, takie że
każde gładkie rozwiązanie równań Seiberga--Wittena spełnia
\[
\|\psi\|_{L^\infty}^2\le
c_0^{-1}\bigl(\tfrac14\|s_g^-\|_{L^\infty}
+C_0\|\eta\|_{L^\infty}\bigr),
\qquad s_g^-:=\max\{0,-s_g\}.
\]
\end{lemat}

\begin{proof}
Korzystamy z normalizacji ustalonej w definicji
\ref{def:tg-sw-eqn} oraz tożsamości z lematu
\ref{lem:tg-sw12-q}. Dla $\Delta=\mathrm d^*\mathrm d$
na funkcjach formuła Weitzenböcka daje
\[
0=\tfrac12\Delta|\psi|^2+|\nabla_A\psi|^2
+\tfrac{s_g}{4}|\psi|^2+\tfrac14|\psi|^4
+\tfrac12\langle\rho(i\eta)\psi,\psi\rangle.
\]
Wynika to z
$\langle\rho(q(\psi))\psi,\psi\rangle=|\psi|^4/2$.
Możemy więc wziąć $c_0=1/4$ i $C_0=C_\rho/2$,
gdzie $\|\rho(i\eta)\|_{\mathrm{op}}\leq C_\rho|\eta|$.
Wartość bezwzględna ostatniego składnika jest
nie większa niż $C_0|\eta||\psi|^2$.

Wybierzmy punkt $x_0$, w którym funkcja $|\psi|^2$
osiąga maksimum. W tym punkcie $\Delta|\psi|^2\geq0$.
Odrzucając również nieujemne $|\nabla_A\psi|^2$
i szacując człon perturbacyjny z dołu, dostajemy
\[
c_0|\psi(x_0)|^4\le
\bigl(\tfrac14s_g^-(x_0)+C_0|\eta(x_0)|\bigr)
|\psi(x_0)|^2.
\]
Jeżeli $\psi(x_0)=0$, spinor jest identycznie zerowy.
W przeciwnym razie dzielimy przez $|\psi(x_0)|^2$
i uzyskujemy żądane oszacowanie globalne.
\end{proof}

\begin{twierdzenie}[Zwartość przestrzeni moduli; twierdzenie analityczne]
\label{thm:tg-sw-compact}
Dla zamkniętej czterorozmaitości $X$, ustalonej struktury
$\mathfrak s$, metryki $g$ i gładkiej perturbacji $\eta$
przestrzeń moduli $\mathcal M_{\mathfrak s}(g,\eta)$
jest zwarta. Każdy ciąg gładkich rozwiązań ma po
zastosowaniu odpowiednich przekształceń cechowania
podciąg zbieżny w topologii $C^\infty$.
\end{twierdzenie}

\begin{proof}[Redukcja do standardowych oszacowań eliptycznych]
Poprzedni lemat daje jednolite ograniczenie
$\|\psi\|_{L^\infty}$. Z drugiego równania wynika
ograniczenie $\|F_A^+\|_{L^2}$.
Ponieważ
$[iF_A/(2\pi)]=c_1(L)\otimes1$, liczba
$\int_X(F_A/i)\wedge(F_A/i)=(2\pi)^2\langle c_1(L)^2,[X]\rangle$
nie zależy od połączenia. Rozkład na części samodualną
i antysamodualną daje zatem także ograniczenie
$\|F_A^-\|_{L^2}$.

Po zapisaniu $A=A_0+a$ wybieramy cechowanie Coulomba
dla części połączeniowej, $\mathrm d^*a=0$, a część harmoniczną $a/i$ przenosimy do
ustalonego zwartego zbioru reprezentantów torusa
$H^1(X;\mathbb R)/(4\pi H^1(X;\mathbb Z))$,
gdzie $H^1(X;\mathbb Z)$ rozumiemy jako obraz w kohomologii
rzeczywistej. Czynnik $4\pi$ wynika z przyjętego działania
$A\mapsto A-2u^{-1}\mathrm du$: forma $u^{-1}\mathrm du/i$
ma okresy należące do $2\pi\mathbb Z$.
Jest to konieczne, gdyż sam warunek Coulomba
nie kontroluje składowej harmonicznej.
Szczegółową konstrukcję takiego cechowania podajemy
w stwierdzeniu \ref{prop:tg-sw12-lattice} i dowodzie
twierdzenia \ref{thm:tg-sw12-compact}.
Istotne jest, że oszacowanie spinora daje w istocie
ograniczenie $F_A^+$ w $L^\infty$, nie tylko w $L^2$.
Dla dowolnego ustalonego $p>4$ z równania
$\mathrm d^+a=F_A^+-F_{A_0}^+$, warunku
$\mathrm d^*a=0$ oraz ograniczenia części harmonicznej
wynika na mocy oszacowania eliptycznego jednostajna
kontrola $\|a\|_{W^{1,p}}$.
Następnie $D_{A_0}\psi=-\rho(a)\psi/2$ daje
ograniczenie $\|\psi\|_{W^{1,p}}$.
Ponieważ $p>4$, przestrzeń $W^{1,p}$ jest algebrą
i zanurza się w $C^0$. Jeśli $a,\psi$ są już ograniczone
w $W^{m,p}$ dla $m\geq1$, to $q(\psi)$ oraz
$\rho(a)\psi$ są ograniczone w tej samej przestrzeni.
Z oszacowań eliptycznych dla
$\mathrm d^*\oplus\mathrm d^+$ i $D_{A_0}$
otrzymujemy ograniczenia w $W^{m+1,p}$.
Indukcja daje wszystkie rzędy regularności bez
kolejnych zmian cechowania. Pełne oszacowania
i wybór podciągu opisano w twierdzeniu
\ref{thm:tg-sw12-compact}.
Na zwartej rozmaitości osadzenia Sobolewa--Rellicha
i argument przekątniowy dają podciąg zbieżny w każdym
$C^m$; granica spełnia równania.
Przyjmujemy standardowe globalne oszacowania eliptyczne
i twierdzenia o osadzeniach Sobolewa, a nie samo
wnioskowanie o kolejnych rzędach regularności.
\end{proof}

\begin{uwaga}[Dlaczego nie ma zjawiska bubblingu?]
Jednolita kontrola spinora oraz abelowy charakter
połączenia pozwalają kontrolować pełną krzywiznę
i jej pochodne po ustaleniu cechowania. Nie pojawiają
się osobne koncentracje krzywizny odpowiadające
nieabelowym instantonom, dla których zwykle trzeba
uzupełniać przestrzeń moduli obiektami osobliwymi.
\end{uwaga}

\section{Transwersalność, orientacje i niezmienniczość}
\label{sec:tg-sw-transversality}

\begin{definicja}[Rozwiązanie regularne]
Nieredukowalne rozwiązanie $(A,\psi)$ nazywamy
\emph{regularnym}, jeżeli $H^2_{A,\psi}=0$.
W takim przypadku $H^1_{A,\psi}$ jest przestrzenią
styczną do lokalnej przestrzeni moduli i ma wymiar
$d(\mathfrak s)$.
\end{definicja}

\begin{twierdzenie}[Regularność generyczna; twierdzenie analityczne]
\label{thm:tg-sw-generic}
Niech $X$ będzie zamkniętą, spójną, zorientowaną
gładką czterorozmaitością z $b_2^+(X)>0$,
a $\mathfrak s$ ustaloną strukturą $\mathrm{Spin}^c$.
Dla odpowiedniego rezydualnego zbioru gładkich perturbacji
samodualnych $\eta$ nie występują rozwiązania redukowalne,
a wszystkie rozwiązania nieredukowalne są regularne.
Wówczas $\mathcal M_{\mathfrak s}(g,\eta)$ jest zwartą,
gładką rozmaitością wymiaru $d(\mathfrak s)$
(lub jest pusta, gdy $d(\mathfrak s)<0$).
\end{twierdzenie}

\begin{proof}[Mechanizm dowodu]
Rozważamy przestrzeń rozwiązań równań ze zmienną
perturbacją $\eta$ i odwzorowanie rzutowania na przestrzeń
perturbacji. Na nieredukowalnej części przestrzeni konfiguracji
działanie grupy cechowania jest wolne. Różniczka równań
łącznie z kierunkiem perturbacji ma składnik
$\delta\eta\mapsto-i\delta\eta$ w równaniu krzywizny.
Surjektywność uniwersalnej linearyzacji wynika
z lematu \ref{lem:tg-sw12-surj}: wariacje perturbacji
zerują składową krzywiznową każdego elementu
ortogonalnego do obrazu, a wariacje połączenia
i jednoznaczne przedłużenie dla operatora Diraca
zerują pozostałą składową.
Po ustaleniu cechowania linearyzacja w samych
konfiguracjach jest operatorem Fredholma $T$.
Niech $P$ oznacza linearyzację w perturbacjach.
Uniwersalna linearyzacja $(T,P)$ jest surjektywna,
więc na dopełnieniach hilbertowskich jej jądro
jest dopełnione. Twierdzenie o funkcji uwikłanej
daje gładką uniwersalną przestrzeń rozwiązań.
To rzutowanie z tej przestrzeni na perturbacje,
a nie uniwersalna linearyzacja z nieskończenie
wymiarowym parametrem, jest odwzorowaniem Fredholma.
Istotnie, na przestrzeni stycznej
$\ker(T,P)$ jego różniczka $(v,h)\mapsto h$
ma jądro $\ker T$, a obraz jest jądrem
surjektywnego odwzorowania
$h\mapsto [Ph]\in\operatorname{coker}T$.
Stąd jego cokernel jest izomorficzny z
$\operatorname{coker}T$, a indeks wynosi
$d(\mathfrak s)$. Wartość regularna rzutowania
odpowiada dokładnie surjektywności $T$.
Twierdzenie Sarda--Smale'a, na rozdzielnych
dopełnieniach Sobolewa i przy klasie różniczkowalności
większej od indeksu, daje zbiór rezydualny
wartości regularnych. Przejście do perturbacji
gładkich wykorzystuje regularność eliptyczną
i standardową wersję generyczności w topologii
$C^\infty$; sam fakt gęstości funkcji gładkich
w normie Sobolewa nie wystarcza do tego przejścia.

Rozwiązanie redukowalne wymaga spełnienia równania w
$\mathcal H_g^+$, które narzuca $b_2^+(X)$ niezależnych
warunków rzeczywistych. Ponieważ $b_2^+>0$, także ten
zbiór perturbacji można generycznie ominąć.
Zwartość wynika z twierdzenia \ref{thm:tg-sw-compact}.
Korzystamy tu z twierdzenia o gładkim ilorazie
nieredukowalnych konfiguracji przez grupę cechowania,
twierdzenia Sarda--Smale'a i eliptycznej regularności.
Dalsze szczegóły linearyzacji i surjektywności
podano w twierdzeniu \ref{thm:tg-sw12-regular}.
\end{proof}

Orientację przestrzeni moduli określa orientacja
\emph{homologiczna}, tj. orientacja rzeczywistej
przestrzeni $H^1(X;\mathbb R)\oplus H^2_+(X;\mathbb R)$.
Jeżeli $b_1=0$, jest to po prostu wybór orientacji
$H^2_+(X;\mathbb R)$. Część zespolona operatora Diraca ma orientację
kanoniczną. Dla części rzeczywistej jądro
$\mathrm d^*\oplus\mathrm d^+$ jest równe
$\mathcal H^1_g$, a cokernel
$\mathcal H^0_g\oplus\mathcal H^+_g$.
Orientujemy $\mathcal H^0_g$ przez funkcję stałą $1$.
Orientacja homologiczna określa więc orientację
linii wyznacznikowej części rzeczywistej.
Homotopia usuwająca człony rzędu zero operatora
deformacyjnego przenosi tę orientację na pełną
linię wyznacznikową. Korzystamy ze standardowej
teorii wiązek wyznacznikowych, opisanej w
stwierdzeniu \ref{prop:tg-sw12-hom-orient};
przy regularnym rozwiązaniu linia ta jest
najwyższą potęgą zewnętrzną przestrzeni stycznej
do przestrzeni moduli.

\begin{stwierdzenie}[Niezmienniczość przy braku redukowalnych]
\label{prop:tg-sw-cobordism}
Niech $b_2^+(X)>1$ oraz
$d(\mathfrak s)=0$. Ustalmy orientację homologiczną
$H^1(X;\mathbb R)\oplus H^2_+(X;\mathbb R)$.
Nie wymagamy $b_1(X)=0$, zgodnie z definicją
niezmiennika w rozdziale \ref{ch:tg-sw}.
Dla dowolnych dwóch generycznych perturbacji i metryk
$(g_0,\eta_0)$, $(g_1,\eta_1)$ otrzymujemy
$SW_{X,g_0,\eta_0}(\mathfrak s)
=SW_{X,g_1,\eta_1}(\mathfrak s)$.
\end{stwierdzenie}

\begin{proof}[Jednowymiarowy kobordyzm]
Łączymy dane generyczne gładką ścieżką
$t\mapsto(g_t,\eta_t)$. Warunek istnienia rozwiązania
redukowalnego ma kowymiar $b_2^+(X)>1$,
więc generyczna ścieżka jednoparametrowa go omija.
Parametryczna wersja twierdzenia Sarda--Smale'a
względem końców ścieżki daje gładką przestrzeń
moduli wymiaru $d(\mathfrak s)+1=1$,
zorientowaną przez orientację linii wyznacznikowej
operatora deformacyjnego i dodatnią orientację
parametru $t$. Zwartość wynika z parametrycznej
wersji oszacowań z twierdzenia \ref{thm:tg-sw-compact}:
dla ciągu parametrów wybieramy podciąg
$t_j\to t_\infty$; w gładkiej rodzinie na
zwartym przedziale metryki są jednostajnie
równoważne, a stałe oszacowań eliptycznych
można dobrać lokalnie jednostajnie wokół
$t_\infty$. Analogicznie kontrolujemy rzut
harmoniczny i kratę cechowania. Argument
przekątniowy prowadzi do rozwiązania dla
danych $(g_{t_\infty},\eta_{t_\infty})$.
Korzystamy w tym kroku z parametrycznej
regularności i standardowej teorii eliptycznej.
Jej brzeg stanowi rozłączną sumę
$\mathcal M_{\mathfrak s}(g_1,\eta_1)$ i przeciwnego
zbioru $\mathcal M_{\mathfrak s}(g_0,\eta_0)$.
Algebraiczna liczba punktów brzegowych zwartej,
zorientowanej jednowymiarowej rozmaitości wynosi zero;
stąd równość liczebności ze znakami.
Wykorzystane istnienie regularnej parametrycznej
przestrzeni moduli jest częścią teorii Fredholma.
\end{proof}

\section{Rozwiązania redukowalne, ściany i komory}
\label{sec:tg-sw-chambers}

Niech $\mathcal H_g^+\subset\Omega^2_+(X)$ oznacza
przestrzeń samodualnych form harmonicznych, a
$\Pi_g^+$ --- rzut $L^2$ na tę przestrzeń.
Dla każdej struktury $\mathfrak s$ z wiązką $L$
ustalmy gładkie połączenie $A_0$ na $L$.

\begin{stwierdzenie}[Równanie ściany]
\label{prop:tg-sw-wall}
Dla ustalonej metryki $g$ i samodualnej perturbacji
$\eta$ istnieje rozwiązanie redukowalne równań
Seiberga--Wittena wtedy i tylko wtedy, gdy
\[
\Pi_g^+\bigl(F_{A_0}/i-\eta\bigr)=0.
\]
Warunek ten nie zależy od wyboru $A_0$. Zbiór perturbacji
spełniających to równanie jest afiniczną podprzestrzenią
o kowymiarze $b_2^+(X)$ w przestrzeni perturbacji.
\end{stwierdzenie}

\begin{proof}
Dla konfiguracji redukowalnej mamy $\psi=0$, więc
$F_A^+=i\eta$. Każde inne połączenie na $L$ ma postać
$A=A_0+ia$ z rzeczywistą formą $a$, a jego krzywizna
wynosi $F_A=F_{A_0}+i\,\mathrm da$.
Szukamy zatem rozwiązania równania
$\mathrm d^+a=\eta-(F_{A_0}/i)^+$.
Na zamkniętej czterorozmaitości operator
$\mathrm d^+:\Omega^1\to\Omega^2_+$ ma domknięty obraz,
a jego dopełnieniem ortogonalnym jest dokładnie
$\mathcal H_g^+$. Jest to następstwo teorii Hodge'a
i eliptyczności operatora $\mathrm d^*\oplus\mathrm d^+$.
Równanie ma więc rozwiązanie dokładnie wtedy, gdy
prawa strona jest ortogonalna do $\mathcal H_g^+$,
czyli gdy zachodzi warunek z tezy.
Zmiana $A_0$ dodaje do krzywizny formę dokładną,
której rzut na $\mathcal H_g^+$ jest zerowy.
Ponieważ rzut $\eta\mapsto\Pi_g^+\eta$ jest surjektywny,
ściana ma kowymiar $\dim\mathcal H_g^+=b_2^+(X)$.
\end{proof}

Dla $b_2^+=1$ ustalmy jedną składową dodatniego stożka
$\{h\in H^2(X;\mathbb R):h^2>0\}$ i niech
$\omega_g$ będzie unormowaną harmoniczną formą
samodualną reprezentującą wybraną orientację $H_g^+$.
Z konwencji krzywizny w definicji
\ref{def:tg-sw-eqn} wynika
$[F_{A_0}/i]=-2\pi c_1(L)\otimes1$.
Wprowadzamy rzeczywistą funkcję
\[
\mathfrak w_{\mathfrak s}(g,\eta)
:=-2\pi\langle c_1(L),[\omega_g]\rangle
-\int_X\eta\wedge\omega_g.
\]
Warunek $\mathfrak w_{\mathfrak s}(g,\eta)=0$ jest dokładnie
równaniem ściany. Komory odpowiadają znakom
$\mathfrak w_{\mathfrak s}>0$ i $\mathfrak w_{\mathfrak s}<0$
po ustaleniu wybranej składowej dodatniego stożka.

\begin{center}
\begin{tikzpicture}[scale=.93,>=Stealth]
  \fill[tgBlue!9] (-4,-1.25) rectangle (0,1.5);
  \fill[tgCoral!10] (0,-1.25) rectangle (4,1.5);
  \draw[tgNavy!55,->] (-4.2,0)--(4.25,0) node[right,font=\scriptsize] {parametr};
  \draw[tgPurple,line width=2.2pt] (0,-1.32)--(0,1.52);
  \node[font=\small,text=tgBlue] at (-2.4,1.02) {$\mathfrak w<0$};
  \node[font=\small,text=tgCoral] at (2.35,1.02) {$\mathfrak w>0$};
  \node[font=\small,text=tgPurple,rotate=90] at (.27,-.63) {ściana redukowalnych};
  \draw[tgTeal,-{Stealth[length=2.9mm]},line width=1.8pt]
     (-2.7,-.48)--(2.6,-.48);
  \node[font=\scriptsize,text=tgNavy] at (-2.7,-.75) {$(g_0,\eta_0)$};
  \node[font=\scriptsize,text=tgNavy] at (2.6,-.75) {$(g_1,\eta_1)$};
  \fill[tgTeal] (-2.7,-.48) circle (.065) (2.6,-.48) circle (.065);
\end{tikzpicture}

{\footnotesize Schemat dwóch komór przy $b_2^+=1$; ścieżka poprzeczna
przecina ścianę w jednym punkcie.}
\end{center}

\begin{lemat}[Poprzeczne przecięcie ściany przez prostą perturbacji]
\label{lem:tg-sw-affine-wall}
Niech $X$ będzie zamkniętą, spójną, zorientowaną,
gładką czterorozmaitością z $b_2^+(X)=1$,
$g$ ustaloną metryką, a $\mathfrak s$
strukturą $\mathrm{Spin}^c$ z wiązką wyznacznikową $L$.
Niech $\omega_g$ będzie harmoniczną formą samodualną
wyznaczającą ustaloną orientację $H_g^+$, znormalizowaną
tak, aby $\int_X\omega_g\wedge\omega_g=1$.
Dla dowolnych ustalonych gładkich rzeczywistych form
samodualnych $\eta_0,\theta$ połóżmy
$\eta_t=\eta_0+t\theta$, $t\in\mathbb R$, i oznaczmy
$a=\int_X\theta\wedge\omega_g$.
Jeżeli $a\ne0$, to rodzina $(g,\eta_t)$ przecina
ścianę redukowalnych dokładnie raz, poprzecznie,
w punkcie
\[
t_*=\frac{-2\pi\langle c_1(L),[\omega_g]\rangle
-\int_X\eta_0\wedge\omega_g}{a}.
\]
Jeżeli $a=0$, cała rodzina albo leży w ścianie,
albo jej nie przecina.
\end{lemat}

\begin{proof}
Z równania ściany
$\mathfrak w_{\mathfrak s}(g,\eta)
=-2\pi\langle c_1(L),[\omega_g]\rangle
-\int_X\eta\wedge\omega_g$
otrzymujemy przez liniowość całki
\[
\mathfrak w_{\mathfrak s}(g,\eta_t)
=\mathfrak w_{\mathfrak s}(g,\eta_0)-ta.
\]
Dla $a\ne0$ funkcja afiniczna ma dokładnie jedno
miejsce zerowe $t_*$, a jej pochodna wynosi $-a\ne0$.
Warunek niezerowej pochodnej jest właśnie
poprzecznością do ściany kowymiaru jeden
wzdłuż ustalonej rodziny. Gdy $a=0$,
funkcja jest stała, co daje ostatnią alternatywę.
Lemat lokalizuje ścianę, ale sam nie dowodzi
formuły zmiany niezmiennika przy przejściu przez nią.
\end{proof}

\begin{przyklad}[Równanie ściany na $\mathbb{CP}^2$]
\label{ex:tg-sw-wall-cp2}
Niech $X=\mathbb{CP}^2$ z orientacją zespoloną,
$H\in H^2(X;\mathbb Z)$ niech oznacza dodatni
generator o $H^2=1$, a
$c_1(L)=3H$. Ustalmy metrykę Fubiniego--Study'ego
i jej harmoniczną formę samodualną
$\omega$ reprezentującą $H$.
Wówczas $\int_X\omega\wedge\omega=1$;
przy tej normalizacji perturbację
wybierzmy w postaci $\eta_t=t\omega$
dla $t\in\mathbb R$.

Ponieważ $\langle3H,[\omega]\rangle=3$,
funkcja definiująca ścianę wynosi
\[
\mathfrak w_{\mathfrak s}(g,\eta_t)
=-2\pi\cdot3-t\int_X\omega\wedge\omega
=-6\pi-t.
\]
Zatem dla $t=-6\pi$ istnieje rozwiązanie
redukowalne równania krzywizny,
a dla $t\ne-6\pi$ takiego rozwiązania nie ma.
Rodzina przecina ścianę poprzecznie,
ponieważ jej pochodna względem $t$
jest równa $-1$.

Mamy $e(X)=3$, $\sigma(X)=1$,
$c_1(L)^2=9$, a więc
$d(\mathfrak s)=(9-6-3)/4=0$.
Dla $t=0$ dodatnia krzywizna skalarna
i równanie Weitzenböcka wykluczają
rozwiązania nieredukowalne,
zaś $-6\pi\ne0$ wyklucza redukowalne.
Zatem niezmiennik w tej komorze
jest zerowy. Twierdzenie o zmianie komory
poniżej, przy jego analitycznych założeniach,
daje w drugiej komorze wartość $\pm1$.
Samo równanie liniowe $-6\pi-t$
nie oblicza znaku tej zmiany.
\end{przyklad}

\begin{twierdzenie}[Zmiana komory w wymiarze zero; wynik analityczny]
\label{thm:tg-sw-wallcross}
Niech $X$ będzie zamkniętą, spójną, zorientowaną,
gładką czterorozmaitością z $b_1(X)=0$ i $b_2^+(X)=1$.
Niech $\mathfrak s$ będzie strukturą $\mathrm{Spin}^c$
o $d(\mathfrak s)=0$. Po wyborze orientacji homologicznej
i spójnej konwencji koorientacji ściany istnieją dwie
liczby $SW_X^+(\mathfrak s)$ i $SW_X^-(\mathfrak s)$,
odpowiadające dwóm komorom, i spełniają one
\[
\bigl|SW_X^+(\mathfrak s)-SW_X^-(\mathfrak s)\bigr|=1.
\]
Znak różnicy zależy od konwencji orientacji oraz tego,
którą komorę oznaczono znakiem $+$.
\end{twierdzenie}

\begin{proof}[Mechanizm lokalny; użycie twierdzenia Kuranishiego]
Wybieramy generyczną jednoparametrową rodzinę danych
przecinającą ścianę poprzecznie dokładnie raz.
Poza ścianą przestrzeń parametryczna ma wymiar jeden.
W punkcie redukowalnym, przy $b_1=0$, jedynym ciągłym
stabilizatorem jest $S^1$. Z rachunku indeksu mamy
\[
0=d(\mathfrak s)=2\operatorname{ind}_{\mathbb C}D_A-2,
\qquad\operatorname{ind}_{\mathbb C}D_A=1.
\]
Dla generycznego przejścia i przy użyciu analitycznej
teorii transwersalności uzyskuje się w punkcie redukowalnym
jednowymiarowe zespolone jądro operatora Diraca oraz
zerowy cokernel. Nie jest to formalny wniosek z samej
równości indeksowej. Lokalna redukcja Lyapunova--Schmidta
i iloraz przez $S^1$ prowadzą do modelu o współrzędnych
$(\tau,z)\in\mathbb R\times\mathbb C$, w którym równanie
ma, z dokładnością do składników wyższego rzędu, postać
$\tau+\varepsilon|z|^2=0$, $\varepsilon\ne0$.
Po podzieleniu przez działanie $z\mapsto e^{i\theta}z$
zostaje półprosta $|z|^2\geq0$, przyłączona do punktu
redukowalnego. Podczas przejścia przez ścianę
pojawia się zatem lub znika jeden punkt nieredukowalny.
Liczenie z orientacjami daje zmianę równą $\pm1$.

Wykorzystany tutaj fakt, że rzeczywisty lokalny model
ma wskazaną postać z niezerowym współczynnikiem,
wymaga analitycznej teorii Kuranishiego i generyczności.
Nie wynika on z samego rachunku indeksu; jest głęboką
częścią twierdzenia o zmianie komory.
\end{proof}

\begin{przyklad}[\texorpdfstring{$\mathbb{CP}^2\#k\overline{\mathbb{CP}}^{\,2}$}{CP2}]
Dla $0\leq k\leq8$ rozważmy
$X_k=\mathbb{CP}^2\#k\overline{\mathbb{CP}}^{\,2}$
z bazą $H,E_1,\ldots,E_k$ taką, że $H^2=1$ i $E_i^2=-1$.
Mamy $e(X_k)=3+k$, $\sigma(X_k)=1-k$, $b_1=0$ i $b_2^+=1$.
Klasa charakterystyczna $c=3H-\sum_{i=1}^kE_i$ spełnia
\[
c^2=9-k=2e(X_k)+3\sigma(X_k),\qquad d(c)=0.
\]
Rozmaitość $X_k$ dopuszcza metrykę o dodatniej krzywiźnie
skalarnej (dla $k>0$ korzystamy z zachowywania tej własności
przez sumę spójną w wymiarze co najmniej trzy).
Przyjmijmy taką metrykę i perturbację $\eta=0$.
Z oszacowania Weitzenböcka wynika, że spinor każdego
rozwiązania musiałby być zerowy. Jednocześnie $c^2=9-k>0$,
a więc samodualny rzut harmoniczny klasy $c$ nie może
być zerowy: gdyby $c^+=0$, mielibyśmy $c^2\leq0$.
Równanie ściany wyklucza zatem także rozwiązania
redukowalne. W tej komorze przestrzeń moduli jest pusta
i niezmiennik wynosi zero. Po przejściu przez ścianę
niezmiennik dla tej samej struktury $\mathrm{Spin}^c$
ma wartość $\pm1$. Nie wolno więc mówić o jednym
niezmienniku Seiberga--Wittena bez określenia komory.
\end{przyklad}

\begin{uwaga}[Dlaczego ma to znaczenie przy wymiernej kontrakcji?]
Wymierna kontrakcja zachowuje $b_2^+$, więc gdy rozmaitość
początkowa ma $b_2^+=1$, rozmaitość wynikowa również ma
$b_2^+=1$. Formuły sklejenia dotyczą wówczas konkretnych
komór zgodnych z klasami użytymi w konstrukcji.
Porównanie z rozdmuchaniami $\mathbb{CP}^2$ musi być
przeprowadzone w tej samej konwencji, a sam brak niezerowych
niezmienników w jednej komorze nie dowodzi niczego
o wszystkich innych komorach.
\end{uwaga}

\section{Zastosowanie: granica dotychczasowych rachunków}
\label{sec:tg-sw-what-next}

Rozpatrzmy raz jeszcze przykład sekcji o samoprzecięciu
$-4$ na powierzchni eliptycznej $E(4)$. Wyjściowa rozmaitość
ma $b_2^+=7$, więc nie występuje zależność niezmiennika od komory.
Przy wymiernej kontrakcji otrzymaliśmy niezmienniki
$e=47$, $\sigma=-31$, a więc $2e+3\sigma=1$.
Dla struktur $\mathrm{Spin}^c$ o oczekiwanym wymiarze zero
koniecznie $c_1(L)^2=1$. Przejście od dwóch ocalałych
klas podstawowych do ich niezerowych wartości wymaga
formuły Fintushela--Sterna oraz danych o przedłużeniach
przez wymierną kulę homologii, nie samego rachunku indeksu.

W przykładach z $b_2^+=1$ dochodzi dodatkowo
wybór komory i weryfikacja poprzeczności rodzin danych.
Dopiero po skompletowaniu: (i) obliczenia grupy fundamentalnej,
(ii) klasyfikacji topologicznej, (iii) zgodnych przedłużeń
$\mathrm{Spin}^c$ i (iv) obliczenia niezmienników w komorze
można poprawnie dowodzić, że skonstruowana rozmaitość
jest homeomorficzna, lecz nie dyfeomorficzna ze standardowym
modelem.

\begin{uwaga}[Zakres dowodów w tej iteracji]
Zostały tutaj bezpośrednio udowodnione: opis stabilizatorów,
linearyzacja równań, równanie ściany, rachunek indeksu
pod warunkiem twierdzenia Atiyaha--Singera, zasada maksimum
dająca oszacowanie spinora oraz konsekwencja kobordyzmu
jednowymiarowego. Przyjęliśmy natomiast analityczne
twierdzenia o przekrojach Coulomba, właściwych własnościach
Fredholma, transwersalności Sarda--Smale'a, pełnej
regularności i zwartości eliptycznej oraz lokalnym modelu
Kuranishiego przy zmianie komory. Oddzielenie tych dwóch
poziomów stanowi część świadomie przyjętego rygoru wykładu.
\end{uwaga}


% ========================================================================
% ITERACJA 12. Umiescic po rozdziale „Analityczne podstawy teorii
% Seiberga--Wittena: przestrzenie moduli i komory”, przed rozdzialem
% o niezmiennikach wezlow. Nie wymaga nowej preambuly.
% ========================================================================
